\section{Forces on the nose}
\label{sec:forces}

\subsection{Axial force from Fluent}

The axial force is integrated over the Fluent wall faces, with $r$ the local radius:
\begin{equation}
F_p = \sum (p_w - p_\infty)\,2\pi r\,\mathrm{d}r, \qquad
F_v = \sum \tau_w\,2\pi r\,\mathrm{d}x, \qquad
C_D = \frac{F_p + F_v}{q\,A_{ref}}, \quad A_{ref} = \pi R^2 = \SI{0.025447}{m^2}.
\end{equation}
The sums reproduce Fluent's own force report to \SI{0.1}{N}. $F_p$ is the pressure drag relative to
$p_\infty$; the net pressure force on the joint of the sealed nose (Table~\ref{tab:lc}) also contains
$(p_\infty - p_{in})\,A_{ref}$.

\begin{table}[H]
\centering
\caption{Axial force at the three Fluent points [N] (smooth wall), and $C_D$ of \zone{wall\_nose}
(pressure + viscous). The coupler carries friction only.}
\label{tab:forces}
\small
\begin{tabular}{@{}lcccccc@{}}
\toprule
 & \multicolumn{3}{c}{A} & \multicolumn{3}{c}{B}\\
\cmidrule(lr){2-4}\cmidrule(l){5-7}
$M$ & 0.95 & 1.20 & \val{pt-M147-M} & 0.95 & 1.20 & \val{pt-M147-M}\\
\midrule
Pressure $F_p$, nose & \val{A-Fp-M095} & \val{A-Fp-M120} & \val{A-Fp-M147} & \val{B-Fp-M095} & \val{B-Fp-M120} & \val{B-Fp-M147}\\
Friction $F_v$, nose & \val{A-Fv-M095} & \val{A-Fv-M120} & \val{A-Fv-M147} & \val{B-Fv-M095} & \val{B-Fv-M120} & \val{B-Fv-M147}\\
Friction, coupler & \val{A-Fc-M095} & \val{A-Fc-M120} & \val{A-Fc-M147} & \val{B-Fc-M095} & \val{B-Fc-M120} & \val{B-Fc-M147}\\
\textbf{Total} & \textbf{\val{A-F-M095}} & \textbf{\val{A-F-M120}} & \textbf{\val{A-F-M147}} & \textbf{\val{B-F-M095}} & \textbf{\val{B-F-M120}} & \textbf{\val{B-F-M147}}\\
\midrule
$C_D$ nose & \val{A-cd-M095} & \val{A-cd-M120} & \val{A-cd-M147} & \val{B-cd-M095} & \val{B-cd-M120} & \val{B-cd-M147}\\
B against A & & & & \val{BvsA-M095}\,\% & \val{BvsA-M120}\,\% & \val{BvsA-M147}\,\%\\
\bottomrule
\end{tabular}
\end{table}

Below $M = 1$ the two noses have the same drag, about \SI{90}{\percent} of it friction. Above $M = 1$
the wave drag dominates, and B's thinner front gives it \val{BlessA-M120} to \val{BlessA-M147}\,\% less.

\subsection{Drag model of the flight}

The flight model (report~2) adds a wave drag from Euler simulations and a flat-plate friction,
\begin{equation}
C_{D,n}(M) = C_{D,w}(M) + C_f\,\frac{S_{wet}}{A_{ref}}, \qquad
C_{D,w}(M) = C_{D,w}(1.1)\,s^2(3 - 2s), \quad s = \frac{M - M_{on}}{1 - M_{on}} \in [0, 1]
\end{equation}
below $M = 1.1$. Report~2 used $M_{on} = 0.8$, which gave 2.6 times Fluent's drag at $M = 0.95$
(\val{A-cdold-M095} against \val{A-cd-M095}). With $M_{on} = 0.94$ the model now
matches Fluent at all three points (Figure~\ref{fig:cd}):
\begin{center}\small
\begin{tabular}{@{}lccc@{}}
\toprule
Fluent / model $- 1$ & $M = 0.95$ & 1.20 & \val{pt-M147-M}\\
\midrule
A & \val{A-dcd-M095}\,\% & \val{A-dcd-M120}\,\% & \val{A-dcd-M147}\,\%\\
B & \val{B-dcd-M095}\,\% & \val{B-dcd-M120}\,\% & \val{B-dcd-M147}\,\%\\
\bottomrule
\end{tabular}
\end{center}
The change raises the apogee by \SI{7}{m} (A) and \SI{1}{m} (B). The rise between $M = 0.95$ and 1.1
is not resolved by Fluent; the model puts it between 0.94 and 1.0.

\begin{figure}[H]
\centering
\includegraphics[width=0.75\linewidth]{fig_cd_mach}
\caption{Nose $C_D$ against Mach along the ascent: flight model with the old and new start of the
drag rise, and Fluent (circles). Smooth wall in both.}
\label{fig:cd}
\end{figure}

\subsection{Axial force along the flight}

\begin{equation}
D_n(t) = q(t)\,A_{ref}\,C_{D,n}\bigl(M(t)\bigr)
\end{equation}
with the friction of a \SI{20}{\micro m} rough wall and the coupler included. It peaks at burnout:
\val{A-nominal-Dnmax}\,N (A) and \val{B-nominal-Dnmax}\,N (B), \val{A-hot-Dnmax} and
\val{B-hot-Dnmax}\,N on the hot day (Figure~\ref{fig:axial}). Fluent, with a smooth wall, gives
\val{A-F-M147} and \val{B-F-M147}\,N at the same point: the two agree within \SI{2}{\percent}.

\begin{figure}[H]
\centering
\includegraphics[width=0.75\linewidth]{fig_axial_force}
\caption{Aerodynamic axial force of the nose along the flight (model) and the Fluent points.}
\label{fig:axial}
\end{figure}

\subsection{Normal force under incidence}

Fluent was run at zero incidence only. Under a small incidence $\alpha$ [rad] slender-body theory gives
the normal force of the nose and its distribution, with $S(x) = \pi r(x)^2$:
\begin{equation}
\frac{\mathrm{d}N}{\mathrm{d}x} = 2\,q\,\alpha\,\frac{\mathrm{d}S}{\mathrm{d}x}, \qquad
N = 2\,q\,A_{ref}\,\alpha, \qquad
x_{CP} = x_S - \frac{V_n}{A_{ref}},
\end{equation}
where $V_n$ is the volume of the nose. Per degree at LC1, $N/\alpha$ = \val{A-nominal-LC1-Nperdeg}\,N
(\val{A-hot-LC1-Nperdeg}\,N hot day) for both designs; $x_{CP}$ = \val{A-xcp}\,mm (A) and
\val{B-xcp}\,mm (B) from the tip. The load is carried mostly by the front half (Figure~\ref{fig:normal});
the coupler cylinder carries none. The bending moment at the joint plane $x_j$ follows as
$M_b = N\,(x_j - x_{CP})$; ST sets the design incidence (Req\_12).

\begin{figure}[H]
\centering
\includegraphics[width=0.75\linewidth]{fig_normal_load}
\caption{Normal load per unit length and per degree of incidence at LC1, slender-body theory; dotted:
centre of pressure.}
\label{fig:normal}
\end{figure}
